\documentclass[11pt]{article}
\usepackage[a4paper,margin=1in]{geometry}
\usepackage{amsmath,amssymb,amsthm,booktabs,hyperref}
\hypersetup{colorlinks=true,linkcolor=blue,urlcolor=blue}
\newtheorem{theorem}{Theorem}
\newtheorem{lemma}{Lemma}
\newcommand{\J}{J_0}
\title{A two-term counterexample to the Bessel-kernel\\
nonnegativity conjecture}
\author{Principia Math in Progress}
\date{12 September 2026}
\begin{document}
\maketitle

\paragraph{Task 1: Resolve MathDB 365454.}
Let \(j_m=j_{0,m}\) be the \(m\)-th positive zero of \(\J\), and put
\[
 K_M^\alpha(x)=
 \sum_{m=1}^M\left(1-\frac{m}{M+1}\right)j_m^\alpha\J(j_mx).
\]
The question asks whether this is nonnegative for every
\(M\ge1\), \(0\le x\le1\), and \(0\le\alpha\le\tfrac12\).
This is Conjecture 1 of Acosta Babb's current author preprint~\cite{preprint}
and of the published paper~\cite{published}, where the kernel is defined
on page 7 and the conjecture appears in the order-zero section on page 11.
We use that intended order-zero case throughout.

\begin{theorem}
The conjecture is false. More precisely,
\[
 K_2^{2/5}(99/100)
 < -\frac{641}{5250000}<-\frac1{10000}.
\]
At the same point, \(K_2^\alpha(99/100)<0\) for every
\(\alpha\ge2/5\), in particular for the whole interval
\([2/5,1/2]\) within the conjectured range. The number of terms \(M=2\)
is minimal for a counterexample.
\end{theorem}

All estimates below are rational inequalities. Numerical approximations
are unnecessary for the proof.

\paragraph{A finite Bessel certificate.}
For rational \(z\), define the rational polynomial value
\[
 P_N(z)=\sum_{k=0}^N\frac{(-1)^k(z^2/4)^k}{(k!)^2}.
\]
For \(0<z<6\), the terms decrease in modulus from index \(3\) onward.
Consequently the alternating-series remainder gives
\begin{equation}\label{alternating}
 P_{13}(z)<\J(z)<P_{12}(z).
\end{equation}
Here is the complete table of finite comparisons used in the proof:
\begin{center}
\begin{tabular}{ccc}
\toprule
\(z\) & polynomial & exact comparison\\
\midrule
\(24048/10000\) & \(P_{13}(z)\) & \(>1/100000\)\\
\(24049/10000\) & \(P_{12}(z)\) & \(<-3/100000\)\\
\(55200/10000\) & \(P_{12}(z)\) & \(<-2/100000\)\\
\(55201/10000\) & \(P_{13}(z)\) & \(>7/1000000\)\\
\(23808/10000\) & \(P_{12}(z)\) & \(<63/5000\)\\
\(54649/10000\) & \(P_{12}(z)\) & \(<-47/2500\)\\
\bottomrule
\end{tabular}
\end{center}
These comparisons can be checked using
\(t_0=1\), \(t_k=-z^2t_{k-1}/(4k^2)\), and
\(P_N=\sum_{k=0}^Nt_k\). The accompanying verifier performs exactly
these operations with integer fractions and records every positive
comparison margin.

\paragraph{Identifying the two zeros.}
Write
\[
 L_1=\frac{24048}{10000},\quad U_1=\frac{24049}{10000},
 \qquad L_2=\frac{55200}{10000},\quad U_2=\frac{55201}{10000}.
\]
The first four rows of the table and~\eqref{alternating} give a zero
in each interval \((L_i,U_i)\). We check their indices, rather than
assuming them from decimal approximations.
The series for \(J_1\) gives, for \(0<x\le5/2\),
\[
 J_1(x)=\frac{x}{2}\sum_{k=0}^\infty
       \frac{(-x^2/4)^k}{k!(k+1)!}
 >\frac{x}{2}\left(1-\frac{x^2}{8}\right)>0.
\]
Thus \(\J'=-J_1<0\) there, and the zero in \((L_1,U_1)\)
is precisely the first positive zero \(j_1\).

To count subsequent zeros, \(u(x)=\sqrt{x}\J(x)\) satisfies
\[
 u''+\left(1+\frac{1}{4x^2}\right)u=0.
\]
If \(a<b\) are any two consecutive zeros with \(a>2\), integration
by parts and the elementary Wirtinger inequality imply
\[
 \frac{\pi^2}{(b-a)^2}\int_a^bu^2
 \le\int_a^b(u')^2
 =\int_a^b\left(1+\frac{1}{4x^2}\right)u^2
 \le\frac{17}{16}\int_a^bu^2.
\]
Hence \(b-a\ge4\pi/\sqrt{17}>2\). Two zero gaps after \(j_1\)
would therefore require a distance greater than \(4\), whereas
\(U_2-L_1<4\). The zero in \((L_2,U_2)\) must be \(j_2\).
This proves the certified enclosures
\begin{equation}\label{zeros}
 L_1<j_1<U_1,\qquad L_2<j_2<U_2.
\end{equation}

\paragraph{Values at the counterexample point.}
The elementary integral representation
\[
 \J(t)=\frac1\pi\int_0^\pi\cos(t\sin\theta)\,d\theta
\]
follows by expanding the cosine and integrating its even powers.
Differentiation under the integral gives \(|\J'(t)|\le1\).
In particular \(\J\) is \(1\)-Lipschitz on the real line.
Put \(x_*=99/100\), \(a_*=23808/10000\), and
\(b_*=54649/10000\). Direct rational interval arithmetic
in~\eqref{zeros} gives
\[
 |x_*j_1-a_*|<\frac1{10000},
 \qquad |x_*j_2-b_*|<\frac1{10000}.
\]
The last two table rows therefore yield
\begin{equation}\label{values}
 0< A:=\J(x_*j_1)<\frac{127}{10000},
 \qquad B:=\J(x_*j_2)<-\frac{187}{10000}.
\end{equation}
The positivity of \(A\) holds because \(0<x_*j_1<j_1\).
Two more finite rational comparisons are
\[
 U_1^2<\left(\frac{10}{7}\right)^5,
 \qquad L_2^2>\left(\frac{49}{25}\right)^5.
\]
Since all quantities are positive, these prove
\[
 j_1^{2/5}<\frac{10}{7},
 \qquad j_2^{2/5}>\frac{49}{25}.
\]
Using the signs in~\eqref{values}, we conclude
\[
\begin{split}
 K_2^{2/5}(x_*)
 &=\frac13\left(2j_1^{2/5}A+j_2^{2/5}B\right)\\
 &<\frac13\left(
 2\frac{10}{7}\frac{127}{10000}
 -\frac{49}{25}\frac{187}{10000}\right)
 =-\frac{641}{5250000}<-\frac1{10000}.
\end{split}
\]
This counterexample has both \(x_*\) and \(\alpha=2/5\) strictly
inside the asserted parameter intervals.

For the stated extension in \(\alpha\), factor
\[
 K_2^\alpha(x_*)=\frac{j_1^\alpha}{3}
       \left(2A+(j_2/j_1)^\alpha B\right).
\]
Because \(j_2/j_1>1\) and \(B<0\), the expression in parentheses
strictly decreases with \(\alpha\). It is already negative at \(2/5\).
Finally \(K_1^\alpha(x)=\tfrac12j_1^\alpha\J(j_1x)\ge0\)
on \(0\le x\le1\), which proves minimality of \(M=2\).
\hfill\(\square\)

\paragraph{Interpretation and verification.}
The conjectured nonnegativity fails; this alone does not disprove the
paper's separate coefficient-existence conjectures.
The exact verifier passes 28 rational checks. Six supplementary checks
at 75 decimal digits give, for example,
\[
 K_2^{2/5}(99/100)\approx-0.00057495277682936770,\qquad
 K_2^{1/2}(99/100)\approx-0.00180586383473877418.
\]
MathDB was still marked open, with no claims in its discussion or
solution fields, during the dated checks saved alongside this note.
No matching full public solution claim was located in the documented
literature and public-catalogue searches.

\paragraph{Task 2: Determine the exact positivity range for \(M=2\).}
Find all real \(\alpha\) for which \(K_2^\alpha(x)\ge0\) on \([0,1]\).
Determine whether the boundary derivative at \(x=1\) detects the
first loss of positivity, or whether an interior point does so first.

\paragraph{Task 3: Analyze the boundary layer for large truncations.}
Obtain an asymptotic expansion of \(K_M^{1/2}(1-y/M)\), including
the dependence on the parity of \(M\), with a uniform error bound
on a specified interval of \(y\). Use it to decide which large
truncations take negative values near the boundary.

\paragraph{Task 4: Construct a positive replacement kernel.}
Find explicit nonnegative weights replacing \(1-m/(M+1)\) that
make the order-zero Bessel kernel nonnegative and retain enough
coefficient control for the applications in the source paper.

\paragraph{Task 5: Resolve the coefficient-existence conjectures directly.}
Prove or refute the endpoint coefficient-existence statement and
its parity-restricted version in the published paper without
assuming nonnegativity of the present kernel. Specify the required
weighted \(L^1\) space in each result.

\begin{thebibliography}{9}
\bibitem{preprint}
R.~L. Acosta Babb,
\emph{A partial converse to the Riemann--Lebesgue lemma for
Bessel--Fourier series of order zero},
arXiv:2410.17681v1, 23 October 2024.
\url{https://arxiv.org/abs/2410.17681}.
\bibitem{published}
R.~L. Acosta Babb,
\emph{On the rate of growth of Bessel--Fourier coefficients for
integrable functions},
J. Math. Anal. Appl. \textbf{559} (2026), 130440.
\url{https://doi.org/10.1016/j.jmaa.2026.130440}.
Published PDF:
\url{https://wrap.warwick.ac.uk/id/eprint/199656/1/1-s2.0-S0022247X26000521-main.pdf}.
\end{thebibliography}
\end{document}
